\documentclass[10pt,a4paper]{article}

% --------------------------------------------------
% PACKAGES
% --------------------------------------------------
\usepackage[utf8]{inputenc}
\usepackage[T1]{fontenc}
\usepackage[french]{babel}

\usepackage{newtxtext}
\usepackage{newtxmath}

\usepackage[
    left=13mm,
    right=13mm,
    top=6mm,
    bottom=6mm
]{geometry}

\usepackage{graphicx}
\usepackage{xcolor}
\usepackage{array}
\usepackage{tabularx}
\usepackage{ragged2e}
\usepackage{enumitem}
\usepackage{fontawesome5}
\usepackage{microtype}
\usepackage{titlesec}
\usepackage{calc}

% --------------------------------------------------
% COULEURS
% --------------------------------------------------
\definecolor{cvburgundy}{RGB}{105,73,84}
\definecolor{cvline}{RGB}{105,73,84}

% --------------------------------------------------
% MISE EN PAGE GENERALE
% --------------------------------------------------
\pagestyle{empty}

\setlength{\parindent}{0pt}
\setlength{\parskip}{0pt}

% --------------------------------------------------
% TITRES DE SECTIONS
% --------------------------------------------------
\titleformat{\section}
  {\normalfont\bfseries\fontsize{12}{14}\selectfont\color{cvburgundy}}
  {}
  {0pt}
  {}
  [\vspace{1pt}\color{cvline}\titlerule]

\titlespacing*{\section}
  {0pt}
  {4pt}
  {2pt}

% --------------------------------------------------
% LISTES
% --------------------------------------------------
\setlist[itemize]{
    leftmargin=1.25em,
    itemsep=0pt,
    topsep=0pt,
    parsep=0pt,
    partopsep=0pt
}

% --------------------------------------------------
% DIMENSIONS
% --------------------------------------------------
\newlength{\datewidth}
\setlength{\datewidth}{2.35cm}

\newlength{\projectgap}
\setlength{\projectgap}{0.25cm}

% --------------------------------------------------
% FORMATION AVEC LOGO
% --------------------------------------------------
\newcommand{\education}[5]{%
    \noindent
    \begin{minipage}[t]{\datewidth}
        \raggedright
        \textbf{#1}
    \end{minipage}%
    \hspace{\projectgap}%
    \begin{minipage}[t]{\dimexpr\linewidth-\datewidth-\projectgap\relax}

        \noindent
        \begin{minipage}[c]{\dimexpr\linewidth-2.50cm\relax}
            \textbf{#2}\\[1pt]
            \textit{#3}\\[1pt]
            #4
        \end{minipage}%
        \hspace{0.15cm}%
        \begin{minipage}[c]{1.20cm}
            \centering
            \ifx&#5&%
            \else
                \includegraphics[
                    height=0.85cm,
                    keepaspectratio
                ]{#5}%
            \fi
        \end{minipage}

    \end{minipage}
    \vspace{2pt}
}

% --------------------------------------------------
% PROJET PRINCIPAL
% --------------------------------------------------
\newcommand{\mainproject}[4]{%
    \noindent
    \begin{minipage}[t]{\datewidth}
        \raggedright
        \textbf{#1}
    \end{minipage}%
    \hspace{\projectgap}%
    \begin{minipage}[t]{\dimexpr\linewidth-\datewidth-\projectgap\relax}
        \noindent
        \textbf{#2}%
        \hfill
        \textit{#3}\\[-10pt]
        #4
    \end{minipage}
    \vspace{3pt}
}

% --------------------------------------------------
% EXPERIENCE
% --------------------------------------------------
\newcommand{\experience}[4]{%
    \noindent
    \begin{minipage}[t]{\datewidth}
        \raggedright
        \textbf{#1}
    \end{minipage}%
    \hspace{\projectgap}%
    \begin{minipage}[t]{\dimexpr\linewidth-\datewidth-\projectgap\relax}
        \noindent
        \textbf{#2}%
        \hfill
        \textit{#3}\\[0pt]
        #4
    \end{minipage}
    \vspace{0pt}
}

% ==================================================
% DOCUMENT
% ==================================================
\begin{document}

% --------------------------------------------------
% EN-TÊTE
% --------------------------------------------------

\noindent
\begin{minipage}[t]{0.76\textwidth}

    {\fontsize{24}{27}\selectfont\bfseries Rayane YAZID}

    \vspace{2pt}

    {\fontsize{10}{12}\selectfont
    Électronique \& Logiciels embarqués
    $\bullet$
    C/C++
    $\bullet$
    Python
    $\bullet$
    ESP32
    $\bullet$
    STM32
    $\bullet$
    FPGA}

    \vspace{7pt}

    \begin{tabular}{@{}p{6.15cm}|@{\hspace{6pt}}p{4.05cm}@{}}

        \faMapMarker* \hspace{2pt} Créteil, 94000
        &
        \faPhone \hspace{2pt} 07 67 14 86 46
        \\[4pt]

        \faEnvelope \hspace{2pt} rayane.yazid.pro@gmail.com
        &
        \faGlobe \hspace{2pt} ryazid.fr
        \\[4pt]

        \faLanguage \hspace{2pt} Anglais : B2
        \hspace{8pt}
        Arabe : B2
        &
        \faCar \hspace{2pt} Permis B

    \end{tabular}

\end{minipage}%
\hfill
\begin{minipage}[t]{0.19\textwidth}

    \centering
    \vspace{0pt}

    \fbox{%
        \includegraphics[
            width=2.3cm,
            height=2.3cm
        ]{pp.jpg}%
    }

\end{minipage}

\vspace{4pt}

% --------------------------------------------------
% PROFIL
% --------------------------------------------------

\section{Profil}

{\small Étudiant en Master E3A à l’Université Paris-Saclay, après une Licence
Sciences pour l’Ingénieur à l’UPEC (6e/76). Je me spécialise en électronique
et systèmes embarqués, avec une expérience en ESP32, STM32, C/C++ et Python.
Je recherche un contrat d’apprentissage ou un stage en systèmes embarqués,
électronique ou robotique.}

% --------------------------------------------------
% FORMATION
% --------------------------------------------------

\section{Formation}

\education
{2026--2028}
{Master 1 -- Électronique, Énergie Électrique et Automatique (E3A)}
{Université Paris-Saclay}
{}
{paris-saclay.png}
\vspace{2pt}
\education
{2023--2026}
{Licence en Sciences pour l’Ingénieur -- classé 6e/76}
{UPEC, Créteil}
{Parcours Électronique / Informatique}
{upec.png}

% --------------------------------------------------
% PROJETS D'INGÉNIERIE
% --------------------------------------------------

\section{Projets d'ingénierie}
\vspace{2pt}
\begingroup\small
\mainproject
{Mai 2026}
{Stage assistant ingénieur -- Bee, diagnostic électronique assisté par IA}
{Phone Addict, Créteil}
{%
\begin{itemize}
    \item \textbf{Instrumentation :} carte de mesure ESP32 (WiFi) : ADC 16~bits ADS1115 en I\textsuperscript{2}C (3 tensions, $\pm$6,144~V), capteur de courant à effet Hall ACS712~5~A (185~mV/A), OLED SSD1306 ; statistiques de session (moyenne, écart-type, tendance) et étude Scilab de l'erreur de quantification de l'ADC.
    \item \textbf{Inférence locale :} serveur Python 3.13 (FastAPI + WebSocket, streaming des réponses) exécutant un LLM Gemma 4B quantifié 4~bits Q4\_K\_M (GGUF, llama-cpp-python, accélération CUDA, RTX 4060 Ti 8~Go, contexte 20\,000 tokens), 100\,\% hors ligne, sans clé API.
    \item \textbf{RAG \& mémoire :} 4 mémoires persistantes (base de connaissances à recherche sémantique sentence-transformers, historique de 600 messages, mémoire permanente, notes de session), 11 compétences auto-détectées, recherche web.
    \item \textbf{Apprentissage \& sortie :} classifieurs Scikit-learn de triage, fine-tuning QLoRA (Unsloth + TRL, $\approx$494 exemples), rapports PDF (fpdf2, PyMuPDF) issus de mesures réelles. Rapport de stage : 16,5/20.
\end{itemize}
}

\mainproject
{Janv. 2026}
{Chariot robotisé -- Châssis mobile articulé}
{UPEC}
{%
\begin{itemize}
    \item \textbf{Motorisation :} 2 moteurs à courant continu EMG30 à encodeurs, alimentés par batterie et commandés en puissance par un pont en H L298 (PWM, sens de rotation) depuis une Arduino Uno.
    \item \textbf{Logiciel \& mécatronique :} firmware C++ de contrôle moteur exploitant les encodeurs (mesure de rotation) ; bras à 3 degrés de liberté intégré sur châssis articulé ; chaîne de puissance, câblage, intégration mécanique.
\end{itemize}
}

\mainproject
{Déc. 2025}
{Bras robotisé -- Commande par gants de données}
{UPEC}
{%
\begin{itemize}
    \item \textbf{Capteurs :} deux gants instrumentés de centrales inertielles MPU6050 (gyroscope + accéléromètre 3 axes, bus I\textsuperscript{2}C) mesurant l'orientation de la main, acquises par un ESP32.
    \item \textbf{Commande :} modèle de cinématique directe convertissant les angles mesurés en consignes articulaires (signaux PWM) pour 4 servomoteurs et une pince pilotée en complément.
    \item \textbf{Logiciel \& mécanique :} firmware C++ d'acquisition et de commande ; liaison filaire à latence de quelques ms (téléopération temps réel) ; interface PC de supervision et de contrôle ; structure conçue et imprimée en 3D.
\end{itemize}
}

% --------------------------------------------------
% EXPÉRIENCE
% --------------------------------------------------

\section{Expérience}
\vspace{2pt}

\experience
{Depuis 2024}
{Concepteur d'objets connectés -- vente en ligne}
{Freelance}
{%
\begin{itemize}
    \item Assistant de bureau connecté (ESP32) destiné à la vente en ligne : firmware C++ en machine à états (yeux animés, heure, météo), OLED SSD1306 en I\textsuperscript{2}C, WiFi configuré au premier démarrage, météo par requêtes HTTP, boîtier conçu sur mesure et imprimé en 3D.
\end{itemize}
}

\vspace{2pt}
\experience
{2023--2026}
{Développeur freelance -- vibe coding \& IA}
{Freelance}
{%
\begin{itemize}
    \item \textbf{Vibe coding \& IA :} 6 sites web (HTML, CSS, JavaScript, Python) en développement assisté par IA avec relecture et validation du code ; analyse fondamentale des LLM : transformer, tokenisation, quantification, contexte, RAG, fine-tuning.
    \item \textbf{Code maintenu :} logiciel de gestion d'atelier Python $\approx$2\,500 lignes (CustomTkinter, SQLite, factures PDF, relances SMTP) ; PWA de suivi de candidatures JavaScript $\approx$1\,500 lignes (Firebase Auth + Firestore en temps réel, 15 utilisateurs).
    \item \textbf{Web, sécurité \& SaaS :} site React 19 / Vite 7 / TypeScript sur Netlify + Cloudflare (fonction serverless, reCAPTCHA v2 vérifié côté serveur avant requête REST Firestore, RGPD, tests Vitest) ; plateforme à abonnement annuel Stripe avec SAS sur mesure (Firebase, 3 clients).
\end{itemize}
}
\endgroup

% --------------------------------------------------
% COMPÉTENCES TECHNIQUES
% --------------------------------------------------

\section{Compétences techniques}
\vspace{1pt}\footnotesize

\textbf{Électronique \& systèmes embarqués :}
C, C++, VHDL (FPGA / CPLD), ESP32, STM32, Basys, Arduino, KiCad, LTspice, Proteus, PSpice,
I2C / SPI / UART / CAN / MQTT, PWM, ADC / DAC, RTOS (FreeRTOS)

\vspace{1pt}

\textbf{Logiciels \& CAO :}
Python (NumPy, SciPy, scikit-learn), MATLAB / Simulink, Scilab, SQL, Vivado, Linux, LabVIEW,
SolidWorks, Fusion 360, Blender 3D, Github, Gitlab, Web (HTML/CSS/JS)


\end{document}